% Applied Software Development Lab -- Lab Handout 6: Moving pictures
% Copyright (c) Michael Lehn. Licensed under CC BY-SA 4.0.
%
% Region coordinates (bp from the top of the page) follow the page layout of
% the original German handout "Laborblatt 6".
\documentclass{aslab}
\sethandout{6}{Moving pictures}
\graphicspath{{../figures/handout06/}{../figures/shared/}}

%% Lists used in this handout
% "A)" items (bold titles, text at +18bp)
\newlist{letters}{enumerate}{3}
\setlist[letters]{label={\Alph*)},leftmargin=18bp,labelsep=*,align=left,
  labelindent=0pt}
% "•" items with the text at +9bp
\newlist{dotlist}{itemize}{3}
\setlist[dotlist]{label={\textbf{\textbullet}},leftmargin=9bp,labelsep=*,
  align=left,labelindent=0pt}
% dash with the blue color of the original (two items on page 4)
\makeatletter
\newcommand*\bluedash{\textcolor{aslabblue}{\aslab@dash}}
\makeatother

%% Input lines "(1) code": monospaced label, code at +4 characters
\newcommand*\inp[1]{\makebox[26.4bp][l]{\ttfamily(#1)}}

%% Python listings with the colors of the original
%%   @..@ orange (line numbers, keywords), ~..~ cyan (built-ins, names),
%%   ^..^ red (numbers), !..! magenta (import), ?..? purple (comments)
\lstdefinestyle{py}{style=aslab,
  moredelim=[is][\color{codeorange}]{@}{@},
  moredelim=[is][\color{codecyan}]{~}{~},
  moredelim=[is][\color{codered}]{^}{^},
  moredelim=[is][\color{codemagenta}]{!}{!},
  moredelim=[is][\color{codepurple}]{?}{?}}
\newcommand*\codexi[1]{\fontsize{11bp}{#1bp}\ttfamily\lineskiplimit=-\maxdimen}

% keep the baseline pitch fixed even next to tall inline code boxes
\lineskiplimit=-\maxdimen

\begin{document}
%% ===================================================================== 1 ==
\maketitleblock

\begin{region}{124.0}
\heading{Learning objectives}
\end{region}

\begin{region}[pitch=13]{138.5}
{\bfseries\color{aslabblue}
\begin{dashlist}
\item Python basics: tuples
\item First steps with PyGame
\item Further steps toward object-oriented programming
\end{dashlist}\par}
\end{region}

\begin{region}{191.5}
\heading{Step 1: Cloning your own repository}
\end{region}

\begin{region}[pitch=14.1]{219.1}
Today we work with Python only. The microcontroller will come back into play
later, though (in the truest sense of the word). So that your code does not
get lost, first clone your repository.
\end{region}

\begin{region}[pitch=12.8]{261.8}
\textbf{As always:}
\begin{dashlist}
\item Open two terminals (left for editing, right for running commands).
\end{dashlist}
\end{region}

\begin{region}{301.2}
\subheading{Left terminal (editing)}
\end{region}

\begin{region}[pitch=13.55]{315.5}
\begin{letters}
\item \textbf{Clone the repository and change the directory}\\
  Clone the repository:\\
  \code{git clone https://gitlab.uni-ulm.de/USERNAME/asp-ulm}
\item \textbf{Change the directory:}\\
  \code{cd asp-ulm/python}
\end{letters}
\end{region}

\begin{region}{395.6}
\subheading{Right terminal (running commands)}
\end{region}

\begin{region}[pitch=12.3]{408.6}
\begin{letters}
\item \textbf{Change the directory}\\
  \code{cd asp-ulm/python}
\end{letters}
\end{region}

\nextpage
%% ===================================================================== 2 ==
\begin{region}{69.4}
\heading{Step 2: Python basics (lists vs.\ tuples)}
\end{region}

\begin{region}[pitch=14]{98.0}
\textbf{Goal:} In the next step, we will use the PyGame module for drawing.
PyGame often uses so-called \textbf{tuples} as a data structure -- for
example, for the size of the window or the position of a graphical
object.\\
Before we really get going with PyGame, let us get familiar with these
\textbf{tuples} and how they differ from \textbf{lists}.
\end{region}

\begin{region}{181.4}
\textbf{Preparation:} In the right window, start a Python shell with the
command \code{python}
\end{region}

\begin{region}{208.8}
\begin{tasks}
\item \textbf{Assignment with several variables}
\end{tasks}
\end{region}

\begin{region}[x=74.7,right=0]{235.3}
Enter the following in the shell:
\end{region}

\begin{region}[x=74.7,right=0,pitch=15.4]{249.9}
\inp{1}\code{x, y = 100, 200}\\
\inp{2}\code{x}\\
\inp{3}\code{y}
\end{region}

\begin{region}[x=74.7,right=0,pitch=14.3]{296.1}
You should observe that the first statement assigns the values \code{100}
and \code{200} to \code{x} and \code{y}, respectively:
\end{region}

\begin{region}[x=74.7,right=0,pitch=13.2]{323.8}
\begin{codeblock}[basicstyle=\codexi{13.2}]
>>> x, y = 100, 200
>>> x
100
>>> y
200
\end{codeblock}
\end{region}

\begin{region}{390.5}
\begin{tasks}[start=2]
\item \textbf{Lists: a mutable data structure}
\end{tasks}
\end{region}

\begin{region}[x=74.7,right=0]{417.3}
Now we create a list \code{pos} and play around with it a little:
\end{region}

\begin{region}[x=74.7,right=0,pitch=15.4]{432.0}
\inp{1}\code{pos = [x,y]}\\
\inp{2}\code{pos[0]}\\
\inp{3}\code{pos[1]}\\
\inp{4}\code{pos}\\
\inp{5}\code{pos.append(42)}\\
\inp{6}\code{pos}
\end{region}

\begin{region}[x=74.7,right=0,pitch=13.9]{522.9}
As you already know from the online Python course, with a list you can:
\begin{dotlist}
\item access individual elements
\item add elements
\end{dotlist}
\end{region}

\begin{region}[x=74.7,right=0]{576.6}
\textbf{In this example, you observe:}
\end{region}

\begin{region}[x=74.7,right=0,pitch=13.2]{590.2}
\begin{codeblock}[basicstyle=\codexi{13.2}]
>>> pos = [x,y]
>>> pos[0]
100
>>> pos[1]
200
>>> pos
[100, 200]
>>> pos.append(42)
>>> pos
[100, 200, 42]
\end{codeblock}
\end{region}

\nextpage
%% ===================================================================== 3 ==
\begin{region}{69.4}
\begin{tasks}[start=3]
\item \textbf{Tuples: an immutable data structure}
\end{tasks}
\end{region}

\begin{region}[x=74.7,right=0,pitch=14.05]{96.9}
Now we create a \textbf{tuple} \code{pos} and observe that you can work with
it in a similar way as with a \textbf{list}.\\
However, a \textbf{tuple} only supports operations that do \emph{not modify}
\code{pos}. This is why \textbf{tuples} are often used to represent fixed
values (e.g., the size of a window or a fixed position).
\end{region}

\begin{region}[x=74.7,right=0]{179.4}
Enter the following in the shell:
\end{region}

\begin{region}[x=74.7,right=0,pitch=15.4]{209.4}
\inp{1}\code{pos = (x,y)}\\
\inp{2}\code{pos[0]}\\
\inp{3}\code{pos[1]}\\
\inp{4}\code{pos}\\
\inp{5}\code{pos[0] = 3}\\
\inp{6}\code{pos.append(42)}
\end{region}

\begin{region}[x=74.7,right=0,pitch=14.0]{313.6}
This is what happened:
\begin{dotlist}
\item With (1), the \textbf{tuple} \code{pos} with two elements was created
  (for this, you use round instead of square brackets).
\item In (2) and (3), the elements were accessed -- this works exactly as
  with a \textbf{list}.
\item In (5), an attempt was made to modify the first element -- this is
  \emph{not allowed} for a \textbf{tuple}.
\item For the same reason, \code{.append()} is not supported either, so a
  tuple cannot be extended.
\end{dotlist}
\end{region}

\begin{region}[x=74.7,right=0]{437.8}
\textbf{Example output:}
\end{region}

\begin{region}[x=74.7,right=0,pitch=13.2]{451.4}
\begin{codeblock}[basicstyle=\codexi{13.2}]
>>> pos = (x,y)
>>> pos[0]
100
>>> pos[1]
200
>>> pos
(100, 200)
>>> pos[0] = 3
Traceback (most recent call last):
  File "<python-input-27>", line 1, in <module>
    pos[0] = 3
    ~~~^^^
TypeError: 'tuple' object does not support item assignment
>>> pos.append(42)
Traceback (most recent call last):
  File "<python-input-28>", line 1, in <module>
    pos.append(42)
    ^^^^^^^^^^
\end{codeblock}
\end{region}

% set in the body font, as in the original
\begin{region}[x=74.7,right=0]{688.8}
AttributeError: \textquotesingle tuple\textquotesingle\ object has no attribute \textquotesingle append\textquotesingle
\end{region}

\begin{region}[pitch=14.35]{716.0}
\textbf{Conclusion:}
\begin{dashlist}
\item \textbf{Tuples} are suitable for \emph{fixed values} that are not
  supposed to change in the program.
\item In \emph{PyGame}, for example, you will often use tuples to pass the
  window size or a fixed position.
\end{dashlist}
\end{region}

\nextpage
%% ===================================================================== 4 ==
\begin{region}{69.4}
\heading{Step 3: First steps with PyGame}
\end{region}

\begin{region}[pitch=13.83]{98.0}
\textbf{Goal:} In this step, we will use the \textbf{SDL} library
(\emph{Simple DirectMedia Layer}).\\
\textbf{SDL} is a library written in C that allows direct access to the
\textbf{graphics card}, the \textbf{sound}, and other \textbf{hardware
components}. This makes it ideal for programming games or other graphical
applications.\\
However, we will use \textbf{SDL not directly}, but through the
\textbf{Python library PyGame}. \textbf{PyGame} is a so-called
\textbf{wrapper} for SDL: it provides us with a convenient Python interface
that lets us use SDL very easily and efficiently.\\
In this step, we will first learn:
\end{region}

\begin{region}[pitch=14.3]{207.9}
\begin{dashlist}
\item how to create a window of a \textbf{specific size} with PyGame, and
\item how to draw a \textbf{graphical object} in it (here: a yellow filled
  circle as the ``ball'').
\end{dashlist}
\end{region}

\begin{region}[pitch=14.3]{237.4}
The main point is to understand how a typical PyGame program is
\textbf{structured}:
\end{region}

\begin{region}[pitch=14.6]{265.0}
\begin{dashlist}
\item[\bluedash] Such a program usually consists of an \textbf{initialization}
  and a \textbf{main loop}, in each iteration of which a \textbf{scene} is
  drawn.
\item[\bluedash] This is the \textbf{basic structure} that we will extend step
  by step later -- for example, \textbf{to move} the ball.
\end{dashlist}
\end{region}

\begin{region}{336.2}
\subheading{Preparation:}
\end{region}

\begin{region}[pitch=13.2]{363.0}
In the left terminal, use the editor of your choice to create a file
\code{breakout.py} with the following content. (The empty lines are
intentional -- we will fill them in bit by bit in the tasks.)
\end{region}

\begin{region}[pitch=12]{416.0}
\begin{codeblock}[style=py,basicstyle=\codexi{12}]
@  1 @!import! pygame
@  2@
@  3@
@  4 @?# setup?
@  5 @pygame.init()
@  6@
@  7 @width, height = ^800^, ^600^
@  8 @screen = pygame.display.set_mode((width, height))
@  9@
@ 10 @ball_c = (^255^, ^255^, ^0^)
@ 11 @ball_x, ball_y = ^400^, ^300^
@ 12@
@ 13 @?# draw stuff?
@ 14@
@ 15 @pygame.draw.circle(screen, ball_c, (ball_x, ball_y), ^5^)
@ 16@
@ 17@
@ 18@
@ 19 @?# done?
@ 20 @pygame.quit()
\end{codeblock}
\end{region}

\begin{region}{669.8}
\subheading{Tasks}
\end{region}

\begin{region}[pitch=13.2]{696.6}
\begin{tasks}
\item Run the program in the right terminal with the command
  \code{python breakout.py}. A new window should appear briefly, but it
  disappears again in a flash.
\end{tasks}
\end{region}

\begin{region}[pitch=13.2]{723.0}
\begin{tasks}[start=2]
\item So far, there has been no explanation of what the individual program
  lines mean. But that is not very interesting at the moment anyway, as
  long as you cannot see anything in the window.\\
  So let us first make sure that the window does not disappear immediately:
\end{tasks}
\end{region}

\begin{region}[x=74.7,right=0,pitch=14.3]{763.4}
\begin{tasks}[label={\bfseries\arabic*.}]
\item In \textbf{line 2}, insert the following line: \code{import time}
\item In \textbf{line 17}, insert the following line: \code{time.sleep(5)}
\end{tasks}
\end{region}

\nextpage
%% ===================================================================== 5 ==
\begin{region}[x=74.7,right=0,pitch=13.4]{68.8}
If you now start the program again with \tcode{python breakout.py}, the
window stays visible for 5~seconds.\\
However, the window is still black and has no visible content.\\
To make the visible result a little more interesting, insert the following
line in \textbf{line 16}:\\
\code{pygame.display.flip()}\\
Run the program once more. Now, during those 5~seconds, you should see a
yellow circle in the middle of the new window.
\end{region}

\begin{region}{176.6}
\subheading{What the program does:}
\end{region}

\begin{region}[pitch=13.5]{189.5}
\begin{dashlist}
\item With \code{import pygame}, the module \code{pygame} is imported.
\end{dashlist}
\end{region}

\begin{region}[pitch=13.9]{203.0}
\begin{dashlist}
\item With \code{pygame.init()}, the necessary preparations are made so
  that you can work with the graphics card, the keyboard, or the sound card
  at all.
\end{dashlist}
\end{region}

\begin{region}[pitch=13.84]{229.7}
\begin{dashlist}
\item With \code{pygame.display.set_mode}, a window is created that you can
  then draw into. The width and height of the window are passed to the
  function -- as a so-called tuple, i.e., a pair of numbers such as
  \code{(800, 600)}.\\
  The function returns a \textbf{variable} that you can later use to draw
  into this window. More precisely, it is a so-called \textbf{object} of
  type \emph{Surface} -- which, in the end, is also nothing but a special
  variable that \textbf{pygame} uses internally to manage the window. For
  us it is enough to know: this variable will later be passed as the first
  argument to drawing functions such as \code{pygame.draw.circle}.
\end{dashlist}
\end{region}

\begin{region}[pitch=13.56]{339.4}
\begin{dashlist}
\item With \code{pygame.draw.circle}, you can draw a circle. The function
  expects four arguments:
  \begin{dashlist}
  \item the variable that represents the window (i.e., \emph{where} to
    draw),
  \item the color of the circle -- here an RGB tuple such as
    \code{(255, 255, 0)} for yellow,
  \item the center of the circle -- as a pair of numbers \code{(x, y)} for
    the position on the screen,
  \item the radius of the circle -- i.e., how large it should be.
  \end{dashlist}
\end{dashlist}
\end{region}

\begin{region}[pitch=13.46]{420.9}
\begin{dashlist}
\item \textbf{pygame} (or rather the underlying \textbf{SDL}) uses a
  technique called \textbf{double buffering}. You can think of it as two
  drawing surfaces: one of them is visible, the other one is ``behind the
  scenes''. You always draw on the hidden surface, and when everything is
  finished, \code{pygame.display.flip()} switches to the new surface. This
  prevents you from watching the program while it draws -- with
  animations, this would otherwise lead to ugly flickering or stuttering.
\end{dashlist}
\end{region}

\begin{region}[pitch=13.55]{501.0}
\begin{dashlist}
\item At the end, the program calls \code{pygame.quit()}. This is not
  strictly necessary, because everything is cleaned up anyway when the
  program terminates. But it is good style to release the resources
  occupied by pygame.
\end{dashlist}
\end{region}

\begin{region}{556.5}
\subheading{To conclude this step:}
\end{region}

\begin{region}[pitch=14.3]{585.5}
Add the file to the repository, create a commit, and push it with the
following commands:
\end{region}

\begin{region}[pitch=13.2]{627.6}
\code{git add breakout.py}\\
\code{git commit -a -m "first steps with pygame"}\\
\code{git push}
\end{region}

\nextpage
%% ===================================================================== 6 ==
\begin{region}{69.4}
\heading{Step 4: Structure of a PyGame program and moving pictures}
\end{region}

\begin{region}[pitch=13.4]{96.9}
\textbf{Goal:}\\
Our PyGame program should not simply run for a certain number of seconds,
but until it is terminated by the user -- either by closing the window (by
clicking the \code{X} at the top right) or by pressing a specific key, for
example \code{Q}.
\end{region}

\begin{region}[pitch=14.17]{177.7}
\textbf{Preparation:}\\
Below you see slightly adapted code for \code{breakout.py} with
\textbf{didactic empty lines} that we will fill in step by step in the
tasks. The main change is the \code{while}\textbf{ loop} in line~16. It runs
as long as the variable \code{running} has the value \code{True}. In this
version, \code{running} is always \code{True} for now. So the statements in
lines 17 to 28 are executed over and over again in an \textbf{infinite
loop}. Another change is that in line~28, after the drawing surface has
been made visible (\code{pygame.display.flip()}), a pause of \textbf{1/30
second} has been added. As a result, the picture is updated \textbf{at
about 30 frames per second (FPS)}. Here is the prepared code:
\end{region}

\begin{region}[pitch=12]{317.8}
\begin{codeblock}[style=py,basicstyle=\codexi{12}]
@  1 @!import! pygame
@  2 @!import! time
@  3@
@  4 @?# setup?
@  5 @pygame.init()
@  6@
@  7 @width, height = ^800^, ^600^
@  8 @screen = pygame.display.set_mode((width, height))
@  9@
@ 10 @ball_c = (^255^, ^255^, ^0^)
@ 11 @ball_x, ball_y = ^300^, ^300^
@ 12@
@ 13@
@ 14 @?# draw stuff?
@ 15 @running = ~True~
@ 16 while@ running:
@ 17@
@ 18@
@ 19@
@ 20@
@ 21@
@ 22@
@ 23@
@ 24 @    pygame.draw.circle(screen, ball_c, (ball_x, ball_y), ^5^)
@ 25@
@ 26@
@ 27 @    pygame.display.flip()
@ 28 @    time.sleep(^1^/^30^)
@ 29@
@ 30 @?# done?
@ 31 @pygame.quit()
\end{codeblock}
\end{region}

\begin{region}{714.5}
\subheading{Tasks}
\end{region}

\begin{region}[pitch=14.25]{728.1}
\begin{tasks}
\item Adapt \code{breakout.py} as shown above and test it in the right
  terminal. The program now runs in an \textbf{infinite loop} and, for the
  moment, can only be stopped by \textbf{aborting} it in the terminal with
  \code{Control-C}.
\end{tasks}
\end{region}

\nextpage
%% ===================================================================== 7 ==
\begin{region}[pitch=13.15]{68.5}
\begin{tasks}[start=2]
\item To make it possible to terminate the program by closing the window,
  insert the following lines. Make sure that the lines are indented so that
  they belong to the while block of the loop in line~16:
\end{tasks}
\end{region}

\begin{region}[x=78.7,right=0]{112.3}
\begin{codeblock}[style=py,basicstyle=\codexi{12}]
@ 17 @    @for@ event @in@ pygame.event.get():
@ 18 @        @if@ event.type == pygame.QUIT:
@ 19 @            running = ~False~
\end{codeblock}
\end{region}

\begin{region}[x=74.7,right=0,pitch=13.93]{156.4}
The function \code{pygame.event.get()} returns a \textbf{list of so-called
events}. An event represents an occurrence -- for example, closing a
window, pressing a key, etc. Test the program before you get a first
\textbf{crash course in ``What is an object?''} afterwards.
\end{region}

\begin{region}[x=74.7,right=0,pitch=13.9]{214.6}
At this point, it can hardly be kept secret any longer: events are
represented as \textbf{objects}. So far, we have mostly avoided the term
object and talked about variables instead. That is not wrong either -- an
object is basically a \textbf{special kind of variable}.
\end{region}

\begin{region}[x=74.7,right=0,pitch=13.9]{272.8}
\textbf{What makes objects special?} An object can consist \textbf{of several
components}. You can think of it as a \textbf{bundle of values} that belong
together. Each of these components can in turn be a variable or another
object. The common term for these components is \textbf{attributes}.
\end{region}

\begin{region}[x=74.7,right=0,pitch=14.27]{328.8}
One such \textbf{attribute} of an event object is, e.g., \code{type}. If an
event represents the \textbf{closing of a window}, this attribute has the
value \code{pygame.QUIT}. In this case, the code above causes
\code{running} to be set to \code{False} -- as a result, the \textbf{outer
while loop} is exited, and the program terminates.
\end{region}

\begin{region}[pitch=13.7]{385.9}
\begin{tasks}[start=3]
\item When a key is pressed on the keyboard, this creates an
  \textbf{event object} whose attribute type has the value
  \code{pygame.KEYDOWN}. Another attribute, \code{key}, stores which key was
  pressed. For example, the value \code{pygame.K_q} represents the key
  \textbf{Q}.\\
  Change the program as follows:
\end{tasks}
\end{region}

\begin{region}[x=78.7,right=0]{444.4}
\begin{codeblock}[style=py,basicstyle=\codexi{12}]
@ 17 @    @for@ event @in@ pygame.event.get():
@ 18 @        @if@ event.type == pygame.QUIT:
@ 19 @            running = ~False~
@ 20 @        @elif@ event.type == pygame.KEYDOWN:
@ 21 @            @if@ event.key == pygame.K_q:
@ 22 @                running = ~False~
\end{codeblock}
\end{region}

\begin{region}[x=74.7,right=0]{524.8}
Now it should also be possible to terminate the program by pressing the
\textbf{Q} key.
\end{region}

\begin{region}[pitch=14.3]{539.1}
\begin{tasks}[start=4]
\item To get some \textbf{movement} into the game, we change the
  \textbf{position of the ball} in every frame. First, we define a
  \textbf{velocity} for the ball in the x and y directions:
\end{tasks}
\end{region}

\begin{region}[x=78.7,right=0]{570.9}
\begin{codeblock}[style=py,basicstyle=\codexi{12}]
@ 12 @ball_vx, ball_vy = -^1^, ^2^
\end{codeblock}
\end{region}

\begin{region}[x=74.7,right=0]{591.4}
The unit of velocity is ``\textbf{pixels per frame}'', because with the lines
\end{region}

\begin{region}[x=78.7,right=0]{608.9}
\begin{codeblock}[style=py,basicstyle=\codexi{12}]
@ 25 @    ball_x += ball_vx
@ 26 @    ball_y += ball_vy
\end{codeblock}
\end{region}

\begin{region}[x=74.7,right=0,pitch=14.2]{641.3}
the \textbf{position} of the ball is updated after about \textbf{1/30
second} each time. Test the program after this change -- the remaining
\textbf{cosmetic flaw} will be fixed next.
\end{region}

\begin{region}[pitch=13.2]{668.8}
\begin{tasks}[start=5]
\item As it moves, the ball still leaves traces. The solution is simple:
  when you draw a new frame, you start by clearing the canvas. Concretely,
  in our case this means that we color it completely black. The RGB value
  for black is (0,0,0), and the following line colors the drawing surface
  black:
\end{tasks}
\end{region}

\begin{region}[x=78.7,right=0]{725.8}
\begin{codeblock}[style=py,basicstyle=\codexi{12}]
@ 23 @    pygame.Surface.fill(screen, (^0^, ^0^, ^0^))
\end{codeblock}
\end{region}

\begin{region}[x=74.7,right=0]{745.0}
The ball should now move without leaving any traces.
\end{region}

\nextpage
%% ===================================================================== 8 ==
\begin{region}[pitch=14.3]{69.4}
\begin{tasks}[start=6]
\item At the moment, the ball still flies off \textbf{into nowhere}. With the
  next change, it should be reflected at the \textbf{edges} of the window.
  From now on, we no longer fill gaps in the program, but \textbf{add new
  lines}.\\
  Add the following lines directly \textbf{after line 26}:
\end{tasks}
\end{region}

\begin{region}[x=78.7,right=0]{129.8}
\begin{codeblock}[style=py,basicstyle=\codexi{12}]
@ 27 @    @if@ ball_x < ^0^ @or@ ball_x > width:
@ 28 @        ball_vx *= -^1^
@ 29 @    @if@ ball_y < ^0^ @or@ ball_y > height:
@ 30 @        ball_vy *= -^1^
\end{codeblock}
\end{region}

\begin{region}[x=74.7,right=0,pitch=14.1]{187.5}
The reflection is not perfect yet, but good enough for now:
\begin{dotlist}
\item If the ball has left the playing field on the \textbf{left or right},
  the velocity in the \textbf{x direction} is reversed.
\item If the ball has left the playing field at the \textbf{top or bottom},
  the velocity in the \textbf{y direction} is reversed.
\end{dotlist}
For a \textbf{perfect reflection}, you would not only have to adjust the
velocity, but also set the \textbf{position} correctly.\\
But that is simple math -- we can add it later.
\end{region}

\begin{region}[pitch=14.11]{300.4}
\begin{tasks}[start=7]
\item The ball flies and bounces -- great! Now we still need a
  \textbf{paddle}, i.e., a \textbf{rectangle}. This is a nice opportunity to
  repeat and vary what we have just learned. In PyGame, a \textbf{paddle} is
  simply a \textbf{rectangle} that is described by:
  \begin{dotlist}
  \item the coordinates of the \textbf{top left} corner and
  \item a \textbf{width} and a \textbf{height}.
  \end{dotlist}
  Of course, you can also give a rectangle a \textbf{color}. To do this, we
  first create a few \textbf{variables}. Add the following lines \textbf{after
  line 12} (analogous to the variables for the ball):
\end{tasks}
\end{region}

\begin{region}[x=78.6,right=0]{416.8}
\begin{codeblock}[style=py,basicstyle=\codexi{12}]
@ 13 @paddle_c = (^0^, ^0^, ^255^)
@ 14 @paddle_x, paddle_y = ^300^, ^570^
@ 15 @paddle_w, paddle_h = ^100^, ^30^
\end{codeblock}
\end{region}

\begin{region}[x=74.6,right=0,pitch=14.3]{465.4}
The paddle should be drawn in every frame. To do this, insert the following
lines \textbf{after line 27} (i.e., \textbf{after the ball has been drawn}):
\end{region}

\begin{region}[x=78.6,right=0]{497.3}
\begin{codeblock}[style=py,basicstyle=\codexi{12}]
@ 28 @    pygame.draw.rect(screen, paddle_c,
 @29@                     (paddle_x, paddle_y, paddle_w, paddle_h))
\end{codeblock}
\end{region}

\begin{region}[pitch=14.27]{532.3}
\begin{tasks}[start=8]
\item Of course, the \textbf{paddle} should \textbf{be movable}. Later, we
  will read the \textbf{serial interface} for this and will then be able to
  change the position by turning the potentiometer. For today, the
  \textbf{left} and \textbf{right} arrow keys will do for steering. We
  already check in the code whether the \textbf{Q} key has been pressed.
  Extend this check as follows:
\end{tasks}
\end{region}

\begin{region}[x=78.7,right=0]{592.7}
\begin{codeblock}[style=py,basicstyle=\codexi{12}]
@ 23 @        @elif@ event.type == pygame.KEYDOWN:
@ 24 @            @if@ event.key == pygame.K_q:
@ 25 @                running = ~False~
@ 26 @            @elif@ event.key == pygame.K_LEFT:
@ 27 @                paddle_x -= ^10^
@ 28 @            @elif@ event.key == pygame.K_RIGHT:
@ 29 @                paddle_x += ^10^
\end{codeblock}
\end{region}

\begin{region}[x=74.7,right=0,pitch=13.2]{698.0}
Test the program after this change. You should notice two things that are
still annoying:
\begin{dotlist}
\item The paddle can be moved out of the window.
\item If you hold down an arrow key, the position does not change
  continuously -- you have to release the key and press it again each time.
\end{dotlist}
\end{region}

\nextpage
%% ===================================================================== 9 ==
\begin{region}[pitch=13.57]{69.4}
\begin{tasks}[start=9]
\item It would be more convenient if, when you \textbf{hold down an arrow
  key}, the paddle kept moving in the corresponding direction until you
  release the key again. To do this, first remove the lines inserted in the
  previous subtask (in my case, these were lines \textbf{26 to 29}).
\end{tasks}
\end{region}

\begin{region}[x=74.7,right=0,pitch=13.8]{124.3}
\textbf{Keyboard events} are only generated when a key is \textbf{pressed} or
\textbf{released} -- \textbf{not} while a key is being \textbf{held down}. If
there were events for that, far too many events would be generated.
\end{region}

\begin{region}[x=74.7,right=0,pitch=14.3]{166.0}
However, you can query \textbf{which keys are currently being held down} and
then check whether the \textbf{arrow keys} are among them.
\end{region}

\begin{region}[x=74.7,right=0,pitch=13.8]{194.6}
To do this, insert the following lines \textbf{after the for loop}. It is
important that these lines belong \textbf{to the while loop}, but \textbf{not
to the for loop} for the events! After all, the lines should also be
executed when there is no event.
\end{region}

\begin{region}[x=74.7,right=0]{236.7}
$\rightarrow$ So pay close attention to the \textbf{indentation depth} (in
my case, 4 characters).
\end{region}

\begin{region}[x=78.7,right=0]{254.3}
\begin{codeblock}[style=py,basicstyle=\codexi{12}]
@ 26 @    keys = pygame.key.get_pressed()
@ 27 @    @if@ keys[pygame.K_LEFT]:
@ 28 @        paddle_x -= ^10^
@ 29 @    @if@ keys[pygame.K_RIGHT]:
@ 30 @        paddle_x += ^10^
\end{codeblock}
\end{region}

\begin{region}[x=74.7,right=0,pitch=13.6]{332.3}
With the function \code{get_pressed()} in line~26, you get a so-called
\textbf{dictionary} in \code{keys}. This data structure is covered in more
detail in the Python bridging course. Here it is enough to know that you
can use it much like a \textbf{list} -- except that the \textbf{index} does
not have to be a number.
\end{region}

\begin{region}[x=74.7,right=0,pitch=13.85]{386.5}
For example, here \code{keys[pygame.K_LEFT]} and \code{keys[pygame.K_RIGHT]}
access entries of \code{keys}. These values are either \code{True} or
\code{False} -- depending on whether the corresponding key is currently
\textbf{pressed or not}.
\end{region}

\begin{region}[pitch=14.3]{428.5}
\begin{tasks}[start=10]
\item Now we also solve the problem that the \textbf{paddle} can currently
  still be moved \textbf{too far to the left} or \textbf{to the right}.\\
  This is easy to do with the following \textbf{fix}, which you should insert
  \textbf{after the position of the paddle may have been changed}:
\end{tasks}
\end{region}

\begin{region}[x=78.7,right=0]{488.9}
\begin{codeblock}[style=py,basicstyle=\codexi{12}]
@ 31 @    @if@ paddle_x < ^0^:
@ 32 @        paddle_x = ^0^
@ 33 @    @if@ paddle_x + paddle_w > width:
@ 34 @        paddle_x = width - paddle_w
\end{codeblock}
\end{region}

\begin{region}{566.9}
\subheading{To conclude this step:}
\end{region}

\begin{region}[pitch=14.4]{583.8}
Add the file to the repository, create a commit, and push it with the
following commands:
\end{region}

\begin{region}[pitch=13.2]{626.0}
\code{git commit -a -m "ball and paddle hack"}\\
\code{git push}
\end{region}

\begin{region}{666.5}
\heading{What comes next}
\end{region}

\begin{region}[pitch=14.3]{682.0}
At this point, you could simply keep working on the game:
\begin{dashlist}
\item For example, make sure that the \textbf{paddle} actually \textbf{makes
  the ball bounce off}.
\item Or that you \textbf{lose the ball} when it flies out of the playing
  field at the bottom.
\end{dashlist}
Unfortunately, though, I will play the \textbf{spoilsport} a little at this
point -- and in the next steps show you instead \textbf{how to organize the
code better}.\\
To do this, we introduce a \textbf{class} \code{Paddle} and a class
\code{Ball}.\\
With this structure, further features can later be implemented \textbf{much
more easily and clearly}.
\end{region}

\nextpage
%% ==================================================================== 10 ==
\begin{region}{69.4}
\heading{Step 5: The class Paddle}
\end{region}

\begin{region}[pitch=14.08]{84.8}
Our \textbf{paddle} consists of \textbf{data} and \textbf{operations} that
manipulate this data:
\begin{dashlist}
\item \textbf{Data} are the variables: \code{paddle_x}, \code{paddle_y},
  \code{paddle_w}, \code{paddle_h}, \code{paddle_c}.
\item \textbf{Operations} are the code that moves the paddle \textbf{to the
  left or right}, as well as the code \textbf{to draw} the paddle.
\end{dashlist}
In the current version, all of this is scattered all over the code. You may
remember from the last lab handout that \textbf{global variables} can be
encapsulated in a \textbf{class}. And you can also put \textbf{functions}
into the class that manipulate these variables.
\end{region}

\begin{region}{195.8}
Applied to the paddle, this could look like this:
\end{region}

\begin{region}[pitch=12]{211.4}
\begin{codeblock}[style=py,basicstyle=\codexi{12}]
@ 14 class@ ~Paddle~:
@ 15 @    c = (^0^, ^0^, ^255^)
@ 16 @    w, h = ^100^, ^30^
@ 17 @    x, y = (width - w) // ^2^, ^570^
@ 18@
@ 19 @    @def@ ~left~():
@ 20 @        Paddle.x -= ^10^
@ 21 @        @if@ Paddle.x < ^0^:
@ 22 @            Paddle.x = ^0^
@ 23@
@ 24 @    @def@ ~right~():
@ 25 @        Paddle.x += ^10^
@ 26 @        @if@ Paddle.x + Paddle.w > width:
@ 27 @            Paddle.x = width - Paddle.w
@ 28@
@ 29 @    @def@ ~draw~(surface):
@ 30 @        pygame.draw.rect(surface, Paddle.c,
@ 31 @                         (Paddle.x, Paddle.y, Paddle.w, Paddle.h))
\end{codeblock}
\end{region}

\begin{region}[pitch=13.3]{429.6}
Replace the lines in which \code{paddle_x}, \code{paddle_y}, \code{paddle_w},
\code{paddle_h}, and \code{paddle_c} were previously initialized with this
code.
\end{region}

\begin{region}[pitch=13.8]{457.0}
Once more, to \textbf{clarify} what this code does:
\begin{dashlist}
\item It initializes the global variables \code{Paddle.c}, \code{Paddle.x},
  \code{Paddle.y}, \code{Paddle.w}, \code{Paddle.h}. You could also say that
  it defines the class attributes \code{c}, \code{x}, \code{y}, \code{w},
  \code{h} of \code{Paddle}.
\item It defines the \textbf{functions} \code{Paddle.left} and
  \code{Paddle.right}, which modify these variables so that the paddle
  moves to the \textbf{left} or \textbf{right}.
\item It defines the \textbf{function} \code{Paddle.draw}, which draws the
  paddle at its current position on a drawing surface.
\item The functions defined in the class could also be called the
  \textbf{class methods} \code{left}, \code{right}, and \code{draw} of
  \code{Paddle}.
\end{dashlist}
\end{region}

\begin{region}[pitch=13.62]{602.9}
\subheading{Task: Make use of it!}
Change the rest of the program so that the \textbf{paddle} is controlled
via the \textbf{class methods}:
\begin{dashlist}
\item if the paddle is to \textbf{move to the left}, call
  \code{Paddle.left()}.
\item if the paddle is to \textbf{move to the right}, call
  \code{Paddle.right()}.
\item And call \code{Paddle.draw(screen)} to draw the paddle on the drawing
  surface \code{screen}.
\end{dashlist}
This does not make the program better (the ball still flies through the
paddle), but it makes it more readable. If everything works:\\
\code{git commit -a -m "paddle class"}\\
\code{git push}
\end{region}

\nextpage
%% ==================================================================== 11 ==
\begin{region}{69.4}
\heading{Step 6: Basics of classes with multiple instances}
\end{region}

\begin{region}[pitch=14.25]{96.9}
For the \textbf{ball}, we could -- as for the paddle -- also define a
\textbf{class}. We now take the opportunity to introduce so-called
\textbf{instantiable classes}. That sounds \textbf{fancy} at first, but it is
quite simple:
\end{region}

\begin{region}[pitch=13.9]{141.9}
It just means that you can create \textbf{several objects} from such a
class.\\
Or, to put it another way:\\
$\rightarrow$ There can be \textbf{several balls} -- but only one
\textbf{paddle}.
\end{region}

\begin{region}[pitch=13.8]{197.2}
We first explore the \textbf{basic idea} in the \textbf{interactive Python
shell}. Start it in the \textbf{right terminal} with the command
\code{python} and first enter the following to define an \textbf{``empty''
class} Ball:
\end{region}

\begin{region}[pitch=13.2]{251.4}
\begin{codeblock}[basicstyle=\codexi{13.2}]
class Ball:
    pass
\end{codeblock}
\end{region}

% "Then:" is set in the monospaced font, as in the original
\begin{region}[pitch=13.2]{295.4}
\tcode{Then:}\\
\inp{1}\code{ball = Ball()}\\
\inp{2}\code{print(type(ball))}\\
\inp{3}\code{ball.x = 42}\\
\inp{4}\code{print(ball.x)}
\end{region}

\begin{region}[pitch=13.7]{374.3}
You should observe:
\begin{dashlist}
\item With (1), a new variable (aka \textbf{object}) \code{ball} was
  defined. As (2) shows, it has the type \code{Ball}.
\item In (3), an \textbf{attribute} \code{x} was added to the object
  \code{ball} and initialized with \code{42}.
\item In (4), this attribute \code{x} is output.
\end{dashlist}
\end{region}

\begin{region}[pitch=13.2]{460.6}
You can also put functions into a class that
\begin{dashlist}
\item attach attributes to the object,
\item manipulate them, or
\item read them.
\end{dashlist}
First, enter the following:
\end{region}

\begin{region}[pitch=13.2]{531.1}
\begin{codeblock}[basicstyle=\codexi{13.2}]
class Ball:
    def init(ich, x):
        ich.x = x
    def print(ich):
        print(f"mein x ist {ich.x}")
\end{codeblock}
\end{region}

\begin{region}[pitch=13.3]{601.2}
Then\\
\inp{1}\code{ball = Ball()}\\
\inp{2}\code{Ball.init(ball, 123)}\\
\inp{3}\code{Ball.print(ball)}
\end{region}

\begin{region}[pitch=13.37]{667.2}
You should observe:
\begin{dashlist}
\item Now \code{ball.x} is initialized with \code{123}.
\item The attribute \code{x} is then output.
\item The advantage: here, the attribute \code{x} is only used and
  \textbf{manipulated} \textbf{inside the functions}.
\end{dashlist}
\end{region}

\nextpage
%% ==================================================================== 12 ==
\begin{region}[pitch=13.3]{69.4}
Now we create a \textbf{second ball} and this time call the functions
\code{init} and \code{print} differently:\\
\inp{1}\code{ball2 = Ball()}\\
\inp{2}\code{ball2.init(456)}\\
\inp{3}\code{ball2.print()}
\end{region}

\begin{region}[pitch=13.15]{148.7}
You should observe:
\begin{dashlist}
\item \code{ball2.init(456)} has the same effect as
  \code{Ball.init(ball2, 456)}.
\item \code{ball2.print()} has the same effect as \code{Ball.print(ball2)}.
\end{dashlist}
\end{region}

\begin{region}[pitch=13.4]{203.0}
\textbf{In fact}, this is exactly how Python \textbf{transforms} it
internally:\\
For the call \code{ball2.print()}, Python automatically turns it into
\code{Ball.print(ball2)} in the background.
\end{region}

\begin{region}[pitch=13.27]{256.2}
Now we try to create a third ball -- but forget to initialize it before we
output it. Enter:\\
\tcode{(1) ball3 = Ball()}\\
\tcode{(2) ball3.print()}
\end{region}

\begin{region}[pitch=13.6]{323.0}
You will see: this results in an \textbf{error}.\\
Because \code{ball3.x} has not been set yet -- this time we did not call
\code{init()}.
\end{region}

\begin{region}[pitch=13.9]{364.0}
\textbf{Wouldn't it be a good idea} if \textbf{initialization} happened
\textbf{automatically} as soon as you create a new ball?\\
$\rightarrow$ Someone actually had this idea already -- and it is built
right into \textbf{Python}!
\end{region}

\begin{region}[pitch=13.2]{418.6}
Enter the following (or use \code{Arrow-Up} to bring back the old code for
the class \code{Ball} and modify it):
\end{region}

\begin{region}[pitch=13.2]{458.4}
\begin{codeblock}[basicstyle=\codexi{13.2}]
class Ball:
    def __init__(ich, x):
        ich.x = x
    def print(ich):
        print(f"mein x ist {ich.x}")
\end{codeblock}
\end{region}

\begin{region}[pitch=13.35]{541.7}
Then enter:\\
\inp{1}\code{ball3 = Ball(321)}\\
\inp{2}\code{ball3.print()}
\end{region}

\begin{region}[pitch=15.0]{595.4}
\textbf{Now you see:}\\
The initialization with \code{__init__} happens automatically as soon as you
create a new object \code{Ball(321)}. So you no longer have to call init
manually -- Python takes care of that \textbf{automatically}.
\end{region}

\begin{region}[pitch=14.3]{652.8}
Usually, you do not use \code{ich} but \code{self} as the \textbf{first
parameter} of the functions (also called \textbf{object methods}); it
specifies the object whose attributes are \textbf{initialized},
\textbf{manipulated}, or \textbf{read}.
\end{region}

\begin{region}[pitch=14.3]{708.8}
However, this is only a \textbf{convention} -- in theory, you could choose
any other name.\\
What is \textbf{not negotiable}, though, is that the function for automatic
initialization (also called \textbf{construction}) must be named
\code{__init__}.
\end{region}

\nextpage
%% ==================================================================== 13 ==
\begin{region}{69.4}
\heading{Step 7: The real class Ball}
\end{region}

\begin{region}[pitch=13.3]{96.9}
The concept of a \textbf{class} from which you can construct \textbf{several
instances} can now be implemented for our game as follows:
\end{region}

\begin{region}[pitch=12]{123.6}
\begin{codeblock}[style=py,basicstyle=\codexi{12}]
@ 10 class@ ~Ball~:
@ 11 @    @def@ ~__init__~(self, x, y, vx, vy):
@ 12 @        self.x = x
@ 13 @        self.y = y
@ 14 @        self.vx = vx
@ 15 @        self.vy = vy
@ 16 @        self.c = (^255^, ^255^, ^0^)
@ 17@
@ 18 @    @def@ ~move~(self):
@ 19 @        self.x += self.vx
@ 20 @        self.y += self.vy
@ 21 @        @if@ self.x < ^0^ @or@ self.x > width:
@ 22 @            self.vx *= -^1^
@ 23 @        @if@ self.y < ^0^ @or@ self.y > height:
@ 24 @            self.vy *= -^1^
@ 25@
@ 26 @    @def@ ~draw~(self, surface):
@ 27 @        pygame.draw.circle(surface, self.c, (self.x, self.y), ^5^)
\end{codeblock}
\end{region}

\begin{region}[pitch=14.3]{340.8}
Use it to replace the lines in which the variables \code{ball_c},
\code{ball_x}, \code{ball_y}, \code{ball_vx}, and \code{ball_vy} were
previously initialized.
\end{region}

\begin{region}{382.1}
\subheading{Brief explanation of the class:}
\end{region}

\begin{region}[pitch=14.2]{395.4}
\begin{dashlist}
\item With the constructor \code{__init__}, you can create a new
  \textbf{Ball object} and specify its \textbf{initial position} and
  \textbf{initial velocity}. Internally, the corresponding
  \textbf{attributes} are attached to the object -- plus an
  \textbf{attribute for the color} with a \textbf{default value}.
\item With the method \code{move}, a Ball object can be \textbf{moved}. At
  the edges, it is \textbf{reflected} as before.
\item With the method \code{draw}, a Ball object can be drawn on a drawing
  surface.
\end{dashlist}
\end{region}

\begin{region}{508.7}
\subheading{Task: Make use of it!}
\end{region}

\begin{region}[pitch=14.75]{522.0}
\begin{dashlist}
\item Every feature you program should also be \textbf{used}. So if we have
  a \textbf{class} from which you can create \textbf{more than one
  instance}, then you should create at least \textbf{two instances}! To do
  this, insert the following line \textbf{before the while loop}:
\end{dashlist}
\end{region}

\begin{region}[x=68.7,right=0]{565.1}
\begin{codeblock}[style=py,basicstyle=\codexi{12}]
@ 49 @balls = [Ball(^400^, ^300^, -^1^, ^2^), Ball(^200^, ^100^, ^4^, ^2^)]
\end{codeblock}
\end{region}

\begin{region}[x=68.7,right=0,pitch=14.2]{579.3}
The \textbf{list} \code{balls} now contains two \textbf{Ball objects} with
different \textbf{initial positions} and \textbf{initial velocities}.
\end{region}

\begin{region}[pitch=15.2]{606.9}
\begin{dashlist}
\item When the balls are \textbf{drawn}, you can \textbf{move them right
  afterwards} as well. To apply this to \textbf{all Ball objects}, you can
  iterate over the elements of \code{balls}:
\end{dashlist}
\end{region}

\begin{region}[x=68.7,right=0,pitch=13.2]{635.7}
\begin{codeblock}[style=py,basicstyle=\codexi{13.2}]
@66 @    @for@ ball @in@ balls:
@67 @        ball.draw(screen)
@68 @        ball.move()
\end{codeblock}
\end{region}

\begin{region}[x=68.7,right=0,pitch=13.9]{676.0}
Use it to replace the line that previously \textbf{drew the ball}. Also
remove the \textbf{old code} that previously moved the ball (i.e., code that
still used \code{ball_x}, \code{ball_y}, \code{ball_vx}, \code{ball_vy}).
\end{region}

\begin{region}[pitch=13.6]{716.6}
\begin{dashlist}
\item If everything works:\\
  \code{git commit -a -m "ball class"}\\
  \code{git push}
\end{dashlist}
\end{region}

\nextpage
%% ==================================================================== 14 ==
\begin{region}{69.4}
\heading{Step 8 (Chill out)}
\end{region}

\begin{region}{98.0}
Extend the program so that a new \textbf{ball} enters the playing field when
the key \code{A} is pressed.
\end{region}

\begin{region}{125.4}
This time you deliberately get \textbf{little help}.
\end{region}

\begin{region}[pitch=13.75]{151.9}
Either you already have an idea or even know how to do it --\\
or you use newfangled stuff like the \textbf{Internet}, which offers things
like \textbf{Google} or \textbf{ChatGPT}.
\end{region}

\begin{region}{205.7}
You just have to ask the right questions, for example:
\end{region}

\begin{region}[pitch=13.2]{231.5}
\begin{dashlist}
\item How do I check in PyGame whether the A key has been pressed?
\item How do I append an object to a list?
\item etc.
\end{dashlist}
\end{region}

\begin{region}[pitch=14.25]{285.7}
The \textbf{new ball} can always have the same \textbf{initial position} or
\textbf{initial velocity} -- or you choose them \textbf{randomly}.\\
And if you think of any other fun additions: \textbf{Go wild!}
\end{region}

\begin{region}{342.7}
\subheading{When you are done -- into the repository with it:}
\end{region}

\begin{region}[pitch=13.2]{370.6}
\code{git commit -a -m "fun with objects"}\\
\code{git push}
\end{region}

\end{document}
